% ---
% id: receiver/def:modes
% kind: def
% title: "Declared mode space"
% summary: >-
%   A declared mode space is a finite, totally ordered set of optical modes,
%   each labelled by a path, a polarization and a reserved temporal or spectral
%   index that takes one declared value. Matrices on the modes follow this
%   order; on the input modes the channel modes come first ([[def:instrument]]).
% prominence: standard
% status: accepted
% uses: []
% mentions: [receiver/component-catalogue/def:catalogue-entries]
% assumes: []
% refs: []
% provenance: newly-derived
% lean:
% ---

The temporal or spectral index of a mode label is reserved: it takes a single
value in this contract, so mode labels carry no temporal semantics. The order compares polarizations last,
so the $H$ and $V$ modes of a path are adjacent.

\begin{definition}\label{def:modes}
A \term[A finite, totally ordered set of optical mode labels together with the Hilbert space whose orthonormal basis they index.]{declared mode space} is a finite totally ordered set $\Modes$, each of
whose elements is an \term[A label recording a spatial path, a polarization and a reserved temporal or spectral index.]{optical mode}, together with the Hilbert space
$\mathbb{C}^{\Modes}$ carrying the orthonormal basis $\{\ket{m}\}_{m\in\Modes}$
indexed in the \tref{declared order}. A mode is a label $m=(p,\pi,\tau)$ recording a
spatial path $p$, a polarization $\pi\in\{H,V\}$ and a temporal/spectral index
$\tau$; the order is lexicographic in $(p,\tau,\pi)$ with $H$ before $V$. The
index $\tau$ is \emph{reserved}: it is carried in the label and in the emitted
artifact, it may take only a single \tref{declared value} in this contract, and no
catalogue entry (Definition~\ref{def:catalogue-entries}) reads or writes it.
For $\Sigma\subseteq\Modes$ we write
$\Emd{\Sigma}=\sum_{m\in\Sigma}\dyad{m}$ for the coordinate projector, and
$\Modes$ is identified with $\{0,\dots,|\Modes|-1\}$ when matrix indices are
needed.
\end{definition}
